\chapter[Band 15: Totale Ordnungen]{Band 15 auf einen Blick}
\noindent\textbf{Totale Ordnungen}\par\medskip
\noindent\emph{Lesart.} Es gilt \(\TotOrd{A,\leq}\),
\(x,y,z\in A\). Bei den binären Operationen \(\min\) und \(\max\)
werden in dieser Übersicht die festen Indizes \({}_{A,\leq}\)
weggelassen.
\begin{BandUebersicht}
\textbf{Totale Ordnung}
Die partielle Ordnung wird um das Vergleichbarkeitsaxiom ergänzt:
\UebFormel{x,y\in A\vdash x\leq y\lor y\leq x.}
\UebQuellen{\FormulaRefAuto{Totale Ordnung}[def];
\FormulaRefAuto{OrderComparability}[ax]}
&
\textit{Minimale und maximale Elemente sind Extremwerte}
Für \(T\subseteq A\):
\UebFormel{\begin{aligned}
&\operatorname{MinEl}_{A,\leq}(m,T)\vdash m=\min_{A,\leq}(T),\\
&\operatorname{MaxEl}_{A,\leq}(m,T)\vdash m=\max_{A,\leq}(T).
\end{aligned}}
\UebQuellen{\FormulaRefAuto{TotalOrderMinimalElementIsMinimum}[thm];
\FormulaRefAuto{TotalOrderMaximalElementIsMaximum}[thm]}\\
\addlinespace[.8ex]
\textbf{Binäre Minimumsoperation}
\UebFormel{\min(x,y)\coloneqq\min_{A,\leq}(\{x,y\}).}
\UebQuellen{\FormulaRefAuto{TotalOrderBinaryMinimum}[def]}
&
\textit{Existenz, Funktionstyp und Infimum}
\UebFormel{\begin{aligned}
&\exists m\in\{x,y\}\,\forall u\in\{x,y\}\,(m\leq u),\\
&\min\colon A\times A\to A,\\
&\min(x,y)=\inf_{A,\leq}(\{x,y\}).
\end{aligned}}
\UebQuellen{\FormulaRefAuto{TotalOrderPairMinimumExistence}[thm];
\FormulaRefAuto{TotalOrderBinaryMinimumFunctionType}[thm];
\FormulaRefAuto{TotalOrderBinaryMinimumAsInfimum}[thm]}\\
\addlinespace[.8ex]
\textbf{Binäre Maximumsoperation}
\UebFormel{\max(x,y)\coloneqq\max_{A,\leq}(\{x,y\}).}
\UebQuellen{\FormulaRefAuto{TotalOrderBinaryMaximum}[def]}
&
\textit{Existenz, Funktionstyp und Supremum}
\UebFormel{\begin{aligned}
&\exists M\in\{x,y\}\,\forall u\in\{x,y\}\,(u\leq M),\\
&\max\colon A\times A\to A,\\
&\max(x,y)=\sup_{A,\leq}(\{x,y\}).
\end{aligned}}
\UebQuellen{\FormulaRefAuto{TotalOrderPairMaximumExistence}[thm];
\FormulaRefAuto{TotalOrderBinaryMaximumFunctionType}[thm];
\FormulaRefAuto{TotalOrderBinaryMaximumAsSupremum}[thm]}\\
\addlinespace[.8ex]
\textbf{Minimum und Maximum als algebraische Operationen}
Die auf allen Paaren definierten Operationen können wiederholt
und in verschiedenen Klammerungen angewandt werden.
&
\textit{Idempotenz, Kommutativität und Assoziativität}
Für \(\star=\min\) und ebenso \(\star=\max\):
\UebFormel{\begin{aligned}
&\star(x,x)=x,\qquad\star(x,y)=\star(y,x),\\
&\star(\star(x,y),z)=\star(x,\star(y,z)).
\end{aligned}}
\UebQuellen{\FormulaRefAuto{TotalOrderBinaryMinimumIdempotent}[thm];
\FormulaRefAuto{TotalOrderBinaryMaximumIdempotent}[thm];
\FormulaRefAuto{TotalOrderBinaryMinimumCommutative}[thm];
\FormulaRefAuto{TotalOrderBinaryMaximumCommutative}[thm];
\FormulaRefAuto{TotalOrderBinaryMinimumAssociative}[thm];
\FormulaRefAuto{TotalOrderBinaryMaximumAssociative}[thm]}\\
\addlinespace[.8ex]
\textbf{Natürliche Ordnung als Beispiel}
Die in Band 10 eingeführte natürliche Ordnung wird mit der
allgemeinen Definition der totalen Ordnung verbunden.
\UebQuellen{\FormulaRefAuto{NaturalOrderOperator}[def]}
&
\textit{Die natürlichen Zahlen sind total geordnet}
\UebFormel{\TotOrd{\mathbb N,\leq_{\mathbb N}}.}
Die einzelnen Ordnungsgesetze aus Band 10 erfüllen gemeinsam
die verlangten Strukturaxiome.
\UebQuellen{\FormulaRefAuto{NaturalNumbersTotalOrder}[thm]}\\
\end{BandUebersicht}
